\scanpage{11}
\mspage{3}
Inversement, de telles données définissent une catégorie $M_0$, satisfaisant $(i)(ii)$.

\uline{NB} Dans les cas que j'ai en tête, $\mathbb N=\mathbb N$

\uline{Ex 1} Catégorie des ens finis tot. ordonnés (ou : des \uline{simplexes} types $\Delta_n=\inserted{[0,n]}$) et des applictions strictement croissantes
% slanted
\margin{$J_n={\unsure{\mathfrak P|\mathscr P}}^{*}(\Delta_n)$\\
$\operatorname{Card} J_n-2^{n+1}-1$}

\uline{Ex 2} Catégorie des \deleted{cubes} \inserted{``\uline{cubes}} épinglés'', \deleted{i.e.} dont leurs objets \unsure{s\textsuperscript{t}|sont} les \deleted{\unsure{systèmes|applications} \unsure{fléchés|fléchées}} cubes combinatoires,
\diagram{\begin{tikzcd}[ampersand replacement=\&, row sep=small,
execute at end picture={\draw[thin] (h104arrbot.south west) -- (h104arrtop.north east);}]
|[alias=h104arrtop]|\deleted{\widetilde A}\arrow[d] \\
|[alias=h104arrbot]|\deleted{A}
\end{tikzcd}}
\deleted{$A,\widetilde A$ ens finis non vides} \deleted{$\widetilde A\unsure{-|=}A$} épinglés par un ordre total sur l'ens des paires d'hyperfaces, \deleted{ou} et le choix d'une hyper\unsure{f.|face}, ou face notée $+$ dans chaque paire.
% slanted
\margin{$J_n={}$ens des sections \unsure{p.tielles|partielles} de $\unsure{\{|(}\,\Delta_{n-1}\times\{\pm1\}\to\Delta_{n-1}\unsure{\}|)}$ avec \unsure{relat\textsuperscript{n}|relat\textsuperscript{o}} d'ordre \uline{opposée} de celle de prolongement\\
$\operatorname{Card} J_n=3^n\deleted{\unsure{4}\,\unsure{2^n}\,\unsure{-2^n}\,\unsure{(2^{n+1}-1)}}$}

T. épinglage du cube induit un épinglage de toute facette, donc ... \deleted{\unsure{I|J|l}}

\uline{Ex. 3} Catégorie \deleted{des} \inserted{des} \uline{hémisphèr\unsure{es}}\deleted{\unsure{iques}} $D_n$, où $J_n=(\Delta_n\times\{\pm1\})\amalg\{d_n\}$, avec $\{\pm1\}$ muni de l'ordre chaotique, $\deleted{\{}d_n\deleted{\}}$ le plus grand élément de $J_n$.
% slanted
\margin{$\operatorname{Card} J_n=2n+1$}

Dans \deleted{tous} ces \inserted{trois} exemples, on a

\textcircled{a} \inserted{Si $n\geq1$} La réalisation top. $|\deleted{\unsure{le}} J_n^*|$ de $J_n^*\inserted{=J_n\setminus\{d_n\}}$ est une \deleted{\unsure{catég|catégorie}}

\deleted{\textcircled{b} \unsure{Sa|La|les} \unsure{réalisation|réalisations}} $(n-1)$ sphère, (donc $|J_n|$ est un $n$-disque)
% archive sheet number: 10
